\ifdefined\isMainDocument
\else
    \documentclass{article}
    \usepackage{fullpage}
    \usepackage{wrapfig}
    \usepackage{lineno}
    \usepackage{amsmath, amsthm, amsfonts, amssymb, amscd}
    \usepackage{bm, bbm, mathrsfs}
    \usepackage{mathtools}
    \usepackage{nicefrac}
    \usepackage{caption}
    \usepackage{float, multirow}
    \usepackage{textcomp, booktabs}
    \usepackage{biblatex}
    \addbibresource{references.bib}
    \usepackage[colorinlistoftodos]{todonotes}
    \usepackage{graphicx}
    \graphicspath{{../main/}{../supp/}}
    \usepackage{setspace}
    \usepackage{hyperref}
    \setstretch{1.5}
    \begin{document}
    \linenumbers
\fi
% Working copy (user, 2026-10-17): scratch.tex is stale; edit this file directly.

\section{Introduction}
% Re-cut 2026-10-21; reassessed 2026-10-23 against the settled Theory, Methods and Results (user:
% "apply all these changes"). Restructured 2026-10-23 (user): the measurement and the two
% questions -> effect, ancestor-centered (Robertson, SC equation) -> effect, gene-centered (Barton
% and Etheridge's bridge; Q = 2; global effective size; multiplicative model) -> linked sites (the
% unit becomes the chromosome segment) -> why the global reduction has been left out (its own
% paragraph, ending on its gap) -> maintenance (MSDB defined here; opens on Li and Berg) -> gap
% and timescale -> this work -> results. Only published results before the last paragraph. The
% two theory paragraphs end on the theory, not on a gap: their gap is the neglect paragraph's.
% Compacted 2026-10-26 (user): restatements of Theory and the Discussion merged or cut; the
% neglect paragraph keeps only the published argument; V_m left to Theory sec:msdb.
% Scaffolding: the voice is the author's to revise.
The additive genetic variance in relative fitness, $V_A$, is the heritable part of the differences in reproductive success among individuals~\cite{crow1958,robertson1961}. In wild populations it can now be measured directly using the pedigree-based animal model, which separates additive genetic variance from the similarity that relatives share through their environment~\cite{kruuk2004}, and when applied to lifetime breeding success it estimates $V_A$ in total fitness rather than in one of its components~\cite{postma2014}. Bonnet et al.~\cite{bonnet2022} meta-analyzed this estimate across 19 long-term monitored populations of birds and mammals and found substantial $V_A$ in most of them, which they interpreted as a sign of ongoing adaptation in phenotypes that share genetic variance with fitness. By the secondary theorem of selection~\cite{robertson1966,price1970}, a non-zero $V_A$ changes the mean breeding value of fitness by $V_A$ each generation, which implies directional selection on the phenotypes that make up fitness without specifying which alleles selection moves nor whether the predicted change is realized, and wild populations often show less evolutionary change than such predictions imply~\cite{merila2001}. At the level of alleles, a variance this large raises two questions: how much neutral diversity it removes, and what maintains it.
% Added 2026-10-26 (user: distinguish selection on phenotypes from selection on alleles, for
% quantitative and population geneticists alike). Robertson 1966 (Anim. Prod. 8:95,
% 10.1017/S0003356100037752), Price 1970 (Nature 227:520, 10.1038/227520a0), Merila et al. 2001
% (Genetica 112-113:199, 10.1023/A:1013391806317): all Crossref-checked 2026-10-26, not yet in
% references.bib. Morrissey et al. 2010 (10.1111/j.1420-9101.2010.02084.x) and the covariance with
% every trait moved to the Discussion only (compaction, 2026-10-26).

Across animal species, the effective size implied by neutral diversity, $\pi = 4N_e\mu$, lies orders of magnitude below the census size~\cite{leffler2012,buffalo2021,lewin2026}, the observation behind Lewontin's paradox~\cite{lewontin1974}. Heritable variance in fitness is one cause of this deficit. Any variance in family size increases drift~\cite{wright1938}, and inheritance prolongs its effects. A favored ancestor passes its advantage to its offspring and, halved by segregation, to theirs, so its descendants remain over-represented for several generations~\cite{robertson1961}. Summed over those generations, the advantage determines the ancestor's expected long-term contribution, and drift at a neutral site is set by how much these contributions vary among ancestors~\cite{wray1990}. Santiago and Caballero~\cite{santiago1995,santiago1998} showed that at long-term equilibrium this reduces effective size, relative to the number of breeding adults $N$, to
\begin{equation}
\frac{N_e}{N} = \frac{1}{1 + Q^2 V_A}, \label{eq:intro_sc}
\end{equation}
where $V_A$ is the long-term equilibrium variance and $Q$ is Robertson's cumulative transmission factor, the ancestor's advantage counted over every generation in which it persists.

From the gene's perspective, Barton and Etheridge~\cite{barton2011} showed that the probability that a gene copy survives is proportional to the reproductive value, the expected long-term contribution, of the individual that carries it, and that selection on unlinked loci in its genetic background can greatly increase the variance of that contribution. A neutral allele therefore shares its carrier's contribution only while it stays associated with the selected alleles that carrier transmitted, and $Q$ is this association summed over the generations it lasts. In population-genetic terms the association is linkage disequilibrium (LD) between the neutral site and the selected loci, a term that applies to unlinked sites as well as to linked ones~\cite{barton2000}. At a site unlinked to every selected locus, the two sites recombine with probability one half, so LD is halved each generation, and summed over all the generations it lasts it comes to twice the carrier's own advantage, $Q = 2$, and selection that removes the selected alleles shortens it further~\cite{santiago1998,charlesworth2012}. Here the individual and gene perspectives are congruent, because the focal allele's only lasting background is its carrier's breeding value. Every site in the carrier's genome shares that value, so the reduction sets a global effective size. Only the gene's perspective extends beyond that, because whatever breaks the association (segregation, recombination or selection) sets how long it persists. Equation~\eqref{eq:intro_sc} holds while $V_A$ is small, for any model of fitness; under multiplicative fitness, on which the animal-model estimate itself rests, the global reduction at larger variances is $e^{-4V_A}$ under random mating and $e^{-8V_A}$ under strict monogamy~\cite{barton2011,weissman2012}.
% Barton and Etheridge 2011 (Genetics 188:953; PMC3176105), abstract: "the probability that a gene
% will survive is proportional to the reproductive value of the individual that carries it";
% "Selection on unlinked loci in the genetic background may greatly increase the variance in
% reproductive value". "Reproductive value" appears here once, as their name for R (CLAUDE.md).

At linked sites the reduction is stronger, because recombination separates a neutral site from a nearby selected locus less often than segregation separates unlinked sites, and it is local, because the chromosome segment rather than the individual carries the association. LD with a selected locus on the same chromosome ends when the two sites recombine or when selection removes the selected allele, so for a tightly linked allele that is purged slowly it persists for many generations and depresses a local effective size that varies along the genome. Under mutation--selection balance this local reduction is background selection~\cite{charlesworth1993,hudson1995}. Santiago and Caballero~\cite{santiago2016} and Buffalo and Kern~\cite{buffaloKern2024} expressed it through the additive variance in fitness, and models of it have been fitted to diversity along the human genome~\cite{murphy2023,buffaloKern2024}, the latter to estimate that variance. Buffalo and Coop~\cite{buffalo2019,buffalo2020} estimated linked selection from temporal change in allele frequencies in \textit{Drosophila}, and across species, linked selection at the strongest published parameters does not account for the deficit in neutral diversity~\cite{buffalo2021}. In each of these studies $V_A$ is inferred from sequence data rather than measured, and the position-by-position predictions need fine-scale recombination maps, whereas the pedigree linkage maps of most wild vertebrates resolve the length of each chromosome far better than the variation along it.
% UNVERIFIED: Murphy et al. 2023 (Murphy, Elyashiv, Amster and Sella, eLife?) not Crossref-checked and not in references.bib.

The global reduction, by contrast, has largely been left out of accounts of the diversity deficit, although at $Q = 2$ it requires only $V_A$ and how long the variance is sustained, with no model of how the variance is maintained and no genomic map. Because LD between unlinked sites is halved every generation, an unlinked association lasts only about two generations, and under mutation--selection balance the unlinked reduction has therefore been taken to be much smaller than the linked reduction and is often omitted~\cite{charlesworth2012,matheson2024}. Yet Matheson and Masel~\cite{matheson2024} found in simulations of a human-like genome that the unlinked effect exceeds the linked effect once the genome-wide deleterious mutation rate is above one. Its size has so far been judged from the variance that deleterious mutation would maintain, not from a variance measured in a population.
% Added 2026-10-23 (user: say why the unlinked term has been left out); own paragraph since the
% restructure. Argument 1 is Matheson and Masel's own account (GBE 16:evae050, 2024): "often
% considered small enough to be ignored", "presumed ... much smaller than that from linked loci
% (Charlesworth 2012)"; they list Elyashiv et al. 2016, Ewing and Jensen 2016, Torres et al. 2018,
% McGee et al. 2022 as models that omit it. The second argument (absorbed into the neutral
% baseline) was OUR reading and was cut here 2026-10-26 (user): only published results before the
% last paragraph; the Discussion makes it, with McVicker et al. 2009.
% UNVERIFIED: Matheson and Masel 2024 not Crossref-checked and not in references.bib.
% AUTHOR: scaffolding, re-voice.

A linked association cannot outlast the selected allele, so how long it persists depends on how fast that allele is purged, and so on what maintains $V_A$~\cite{li2026}. The classical account of what maintains it is mutation--selection--drift balance (MSDB), between recurrent deleterious mutation and purifying selection~\cite{haldane1937,crow1970}. When selection is strong relative to drift, the variance maintained is
\begin{equation}
V_A \approx U_d \bar{s}, \label{eq:intro_mukai}
\end{equation}
where $U_d$ is the diploid genome-wide deleterious mutation rate and $\bar{s}$ the mean heterozygous effect of new mutations on relative fitness, the mean of the distribution of fitness effects (DFE)~\cite{mukai1972,charlesworth2015}. At this balance selection is purifying at every locus, and the gain in mean breeding value for fitness that it makes each generation, $U_d\bar{s}$, is offset by new mutations, so the phenotypes that make up fitness are under directional selection while their breeding values do not change. Charlesworth~\cite{charlesworth2015} compared this model against the genetic variance in fitness components of \textit{Drosophila} and found more than mutation could maintain, pointing to balancing selection instead~\cite[see also][]{johnson2005}, as did mutation-accumulation lines for two of three fitness components~\cite{sharp2018}. For wild vertebrates, Connallon and Czuppon~\cite{connallon2026}, modeling how balancing selection can maintain such variance, noted that deleterious mutation rates appear too low to account for the estimates of Bonnet et al.~\cite{bonnet2022}. That argument, however, was qualitative. Both inputs of Equation~\eqref{eq:intro_mukai} are now measured in vertebrates: $U_d$ from mutation rates estimated by pedigree sequencing~\cite{bergeron2023} and annotated genomes, and $\bar{s}$ from estimates of the DFE~\cite{eyrewalker2007,huber2017,lin2026}.

To our knowledge, pedigree estimates of total fitness $V_A$ in wild vertebrates have not been tested quantitatively against either expectation. Bonnet et al.'s estimates are additive variance in relative fitness, the quantity both theories use, so neither test requires a change of units. A pedigree estimate, however, describes the cohorts and environments sampled, and heritability in wild populations changes with environmental conditions~\cite{charmantier2005,wilson2006,mawass2022}. Both expectations instead concern the long-term variance. Neutral diversity reflects effective size over the roughly $4N_e$ generations of a coalescence, and MSDB describes a variance that holds constant over time, which is the stationarity the diversity test assumes and the maintenance test examines. Each test is therefore a statement about how a few decades of variance in fitness relate to its long-term value.

Here we put the 19 pedigree estimates of Bonnet et al.~\cite{bonnet2022} to both tests, using the multiplicative model for the estimate and the prediction alike. We extend the global reduction of Equation~\eqref{eq:intro_sc} to separate sexes and social monogamy, so that a single $V_A$ estimated across both sexes applies to every mating system~\cite{santiago1995,woolliams1999}, and extend the multiplicative bound, the largest reduction a given $V_A$ can produce, to extra-pair paternity. We add linked sites by averaging the local reduction over position along each chromosome, which requires only its map length, and use measured mutation rates to find the mean selection coefficient needed to maintain each estimate. Finally, we ask how far the long-term variance would have to differ from the pedigree estimate to pass each test. A departure in one test could reflect an unrepresentative study period or processes the model leaves out. A larger long-term variance, however, removes more diversity but is harder to maintain, so any error in the estimate narrows one discrepancy and widens the other, and since a population has only one long-term variance, the two tests together can separate the alternatives.

Our results reveal that the measured $V_A$ is too small to explain the neutral diversity these populations are missing, and too large to be maintained at MSDB. Heritable variance in fitness roughly halves the global effective size in the median population, whereas observed diversity lies hundreds to tens of thousands of times below its neutral expectation, so that even under the multiplicative bound the predicted diversity exceeds the observed by factors of 80 to 31,000; explaining the deficits would take a long-term $V_A$ above one, more than twice the largest estimate. Linkage lowers that requirement, though not below any population's credible interval, and only when the variance is carried by slowly purged alleles, which could not in turn maintain it. Maintenance would require a mean selection coefficient about 18 times the DFE means assigned to these populations. The long-term variances that the two tests require differ by a factor of about 390 on average, and by more than 100 in every population, so no single rescaling of the pedigree estimates satisfies both. The conflict does not depend on the estimates, because the variance that recurrent mutation can maintain at the measured DFEs would remove only a few percent of neutral diversity at unlinked sites, and linkage barely changes this.

\ifdefined\isMainDocument
\else
    \printbibliography
    \end{document}
\fi
